\chapter{开放量子系统}
\label{chap:phys-oqs-open-quantum-systems}

封闭复合系统的总态始终幺正演化，纯态不会凭空变成混态。实验却常常只读取其中一部分自由度：探测原子而不记录逸出的光子，或读取量子比特而不测量它所耦合的振子。丢掉未观测结果以后，系统的相干性可以下降，即使总态仍然完全纯。

这个现象不需要弱耦合，也不需要马尔可夫环境。两个二能级已经足以看清全部逻辑：联合纯态的密度算符取偏迹后，环境标签相互正交的交叉项消失，系统态便可能变混。把同一计算写在任意环境基中，就得到 Kraus 算符及其完备关系；而一旦初态含有系统--环境关联，“给任意系统态规定同一个环境态”的制备假设便不再成立。

在开始计算以前，先把态矢量、密度算符和偏迹之间的关系从测量概率重新建立起来。这样做会多用几页，却能避免后面把“总态的部分信息没有读取”误解成“系统真的随机跳到了某个未知纯态”。

\section{从态矢量到密度算符}

量子力学中，归一化态矢量 \(\lvert\psi\rangle\) 对应的测量平均值是
\[
\langle A\rangle_\psi
=\langle\psi\rvert A\lvert\psi\rangle.
\]
利用迹的循环性，可以把它写成
\[
\langle A\rangle_\psi
=\Tr\!\left(\lvert\psi\rangle\langle\psi\rvert A\right).
\]
这提示我们把
\(\rho_\psi=\lvert\psi\rangle\langle\psi\rvert\)
作为纯态的另一种写法。它看似只是记号变化，却能描述实验制备中常见的经典不确定性：若制备装置以概率 \(p_k\) 输出态
\(\lvert\psi_k\rangle\)，所有测量平均值都由
\begin{equation}
\rho=\sum_k p_k\lvert\psi_k\rangle\langle\psi_k\rvert,
\qquad p_k\geq0,\qquad \sum_k p_k=1
\label{eq:oqs-density-mixture}
\end{equation}
给出，因为
\[
\sum_kp_k\langle\psi_k\rvert A\lvert\psi_k\rangle
=\Tr(\rho A).
\]

\begin{definitionplain}[密度算符]
Hilbert 空间上的算符 \(\rho\) 若满足
\[
\rho=\rho^\dagger,
\qquad \langle\phi\rvert\rho\lvert\phi\rangle\geq0,
\qquad \Tr\rho=1,
\]
就称为密度算符。三个条件分别保证概率为实数、非负并且总和为一。
\end{definitionplain}

正算符的本征值都非负。有限维谱定理于是给出
\[
\rho=\sum_j\lambda_j\lvert j\rangle\langle j\rvert,
\qquad \lambda_j\geq0,\qquad \sum_j\lambda_j=1.
\]
所以每个密度算符都能解释成某个正交态系综。这个分解一般不唯一；真正可观测的是 \(\rho\)，不是制备者选择的某一份系综清单。

\begin{proposition}[纯度判据]
若 Hilbert 空间维数为 \(d\)，则
\[
\frac1d\leq\Tr\rho^2\leq1.
\]
等号 \(\Tr\rho^2=1\) 当且仅当 \(\rho\) 是秩一投影；此时状态为纯态。
\end{proposition}

\begin{proof}
在 \(\rho\) 的本征基中，
\(\Tr\rho^2=\sum_j\lambda_j^2\)。由于
\(0\leq\lambda_j\leq1\)，有
\(\sum_j\lambda_j^2\leq\sum_j\lambda_j=1\)。等号要求每个
\(\lambda_j^2=\lambda_j\)，再结合本征值之和为一，只能有一个本征值为一，其余为零。下界由 Cauchy--Schwarz 不等式
\[
1=\left(\sum_j\lambda_j\right)^2
\leq d\sum_j\lambda_j^2
\]
得到；所有本征值都等于 \(1/d\) 时取等号。
\end{proof}

复合系统 \(S+B\) 的 Hilbert 空间是张量积
\(\mathcal H_S\otimes\mathcal H_B\)。乘积态
\(\lvert s\rangle\otimes\lvert b\rangle\)
常简写成 \(\lvert s,b\rangle\)。并非每个联合态都能写成一个乘积。例如
\[
\lvert\Phi^+\rangle
=\frac{\lvert0,0\rangle+\lvert1,1\rangle}{\sqrt2}
\]
若能写成
\((a\lvert0\rangle+b\lvert1\rangle)\otimes
(c\lvert0\rangle+d\lvert1\rangle)\)，
则 \(ac=bd=1/\sqrt2\) 而 \(ad=bc=0\)，这四个条件互相矛盾。因此它是纠缠态。

实验只测系统 \(S\) 时，需要一个系统算符 \(\rho_S\)，使每个只作用在 \(S\) 上的可观测量 \(A\) 都满足
\begin{equation}
\Tr_{SB}\!\left[\rho_{SB}(A\otimes I_B)\right]
=\Tr_S(\rho_S A).
\label{eq:oqs-partial-trace-operational}
\end{equation}
这条对所有 \(A\) 成立的要求定义了偏迹，而不是先凭形式删掉环境指标。

\begin{definitionplain}[偏迹]
取 \(\mathcal H_B\) 的任意正交归一基
\(\{\lvert\alpha\rangle\}\)。对联合算符 \(X\)，定义
\begin{equation}
\Tr_B X
=\sum_\alpha
(I_S\otimes\langle\alpha\rvert)X
(I_S\otimes\lvert\alpha\rangle).
\label{eq:oqs-partial-trace-definition}
\end{equation}
结果是只作用在 \(\mathcal H_S\) 上的算符。
\end{definitionplain}

\begin{proposition}[偏迹与所选环境基无关]
式 \eqnrefs{eq:oqs-partial-trace-definition} 在任何环境正交归一基中给出同一个算符，并且满足操作定义
\eqnrefs{eq:oqs-partial-trace-operational}。
\end{proposition}

\begin{proof}
设另一组基为
\(\lvert\mu\rangle=\sum_\alpha U_{\alpha\mu}\lvert\alpha\rangle\)，
其中 \(U\) 酉。用新基求和得到
\[
\begin{aligned}
\sum_\mu(I\otimes\langle\mu\rvert)X(I\otimes\lvert\mu\rangle)
&=\sum_{\mu,\alpha,\beta}
U_{\alpha\mu}^*U_{\beta\mu}
(I\otimes\langle\alpha\rvert)X(I\otimes\lvert\beta\rangle)\\
&=\sum_\alpha
(I\otimes\langle\alpha\rvert)X(I\otimes\lvert\alpha\rangle),
\end{aligned}
\]
第二行使用
\(\sum_\mu U_{\alpha\mu}^*U_{\beta\mu}=\delta_{\alpha\beta}\)。
再在 \(S\) 的一组基 \(\{\lvert i\rangle\}\) 中展开迹，便有
\[
\begin{aligned}
\Tr_S[(\Tr_BX)A]
&=\sum_{i,\alpha}
\langle i,\alpha\rvert X(A\otimes I)\lvert i,\alpha\rangle\\
&=\Tr_{SB}[X(A\otimes I)].
\end{aligned}
\]
因此偏迹恰好保留所有系统局域测量的统计，而且不依赖计算时使用的环境基。
\end{proof}

三个简单例子可以把不同来源的混合区分开。乘积态
\(\rho_S\otimes\rho_B\) 满足
\(\Tr_B(\rho_S\otimes\rho_B)=\rho_S\)。经典相关态
\[
\rho_{SB}^{\rm cl}
=\tfrac12\lvert00\rangle\langle00\rvert
+\tfrac12\lvert11\rangle\langle11\rvert
\]
与纠缠纯态
\(\lvert\Phi^+\rangle\langle\Phi^+\rvert\)
具有同一个约化态 \(I/2\)，但前者的联合纯度为 \(1/2\)，后者为 \(1\)。只看 \(S\) 无法区分二者；必须测量系统--环境关联。这也说明约化态变混并不等于联合态发生了随机坍缩。

\section{偏迹与约化态}

设系统 $S$ 与环境 $B$ 都是二能级。相互作用把初始乘积态演化为
\begin{equation}
|\Psi\rangle=a|0\rangle_S|0\rangle_B+b|1\rangle_S|1\rangle_B,
\qquad |a|^2+|b|^2=1.
\label{eq:oqs-two-qubit-entangled}
\end{equation}
联合密度算符含有四项，
\begin{align*}
|\Psi\rangle\langle\Psi|
={}&|a|^2|00\rangle\langle00|+|b|^2|11\rangle\langle11|\\
&+ab^*|00\rangle\langle11|+a^*b|11\rangle\langle00|.
\end{align*}
对环境取偏迹就是在一组完备环境基上夹取并求和。交叉项含有
$\langle1|0\rangle_B$ 或 $\langle0|1\rangle_B$，因此消失，留下
\begin{equation}
\rho_S=\Tr_B|\Psi\rangle\langle\Psi|
=|a|^2|0\rangle\langle0|+|b|^2|1\rangle\langle1|.
\label{eq:oqs-two-qubit-reduced}
\end{equation}
系统纯度为 $\Tr\rho_S^2=|a|^4+|b|^4$。只要 $ab\neq0$，它就小于一；当 $|a|=|b|=1/\sqrt2$ 时纯度降到 $1/2$。相干项并没有在联合态中消失，而是储存在系统与环境的关联里；只看系统时，环境基态的正交性把这些项从可观测统计中删掉。

\section{精确约化动力学}
\label{sec:phys-oqs-exact}

把所有未直接观测的自由度统称为环境 $B$。下面的总演化始终保持精确，近似只会在下一章为了得到时间局域方程时出现。

环境的自由哈密顿量记为 $H_B$。系统--环境相互作用写成一般双线性形式
\begin{equation}
H_{SB}=\lambda H_I,
\qquad
H_I=\sum_\mu S_\mu\otimes B_\mu,
\label{eq:phys-oqs-general-system-bath}
\end{equation}
其中 $S_\mu$ 与 $B_\mu$ 分别只作用于系统和环境 Hilbert 空间，$\lambda$ 取为无量纲的记账参数，只标记系统--环境耦合的微扰阶数；$H_I$ 本身保留能量量纲。总哈密顿量为
\begin{equation}
H_{\rm tot}=H_S+H_B+\lambda H_I.
\label{eq:phys-oqs-total-h}
\end{equation}
对任意给定初始总态 $\rho_{\rm tot}(0)$，严格总演化为
\[
\rho_{\rm tot}(t)
=U_{\rm tot}(t,0)\rho_{\rm tot}(0)U_{\rm tot}^\dagger(t,0),
\]
因此系统约化态定义为
\begin{equation}
\rho(t)=\Tr_B\!\left[
U_{\rm tot}(t,0)\rho_{\rm tot}(0)U_{\rm tot}^\dagger(t,0)
\right].
\label{eq:phys-oqs-exact-reduced-state}
\end{equation}
\eqnrefs{eq:phys-oqs-exact-reduced-state} 是开放系统理论的严格起点，不要求弱耦合、马尔可夫、久期或热平衡。

若进一步固定环境初态 $\rho_B$ 并采用无初始相关的制备
\begin{equation}
\rho_{\rm tot}(0)=\rho(0)\otimes\rho_B,
\label{eq:phys-oqs-factorized-initial-state}
\end{equation}
则可把约化演化定义为对任意系统初态都适用的映射
\begin{equation}
\Phi_t[\rho(0)]
=\Tr_B\!\left[
U_{\rm tot}(t,0)(\rho(0)\otimes\rho_B)U_{\rm tot}^\dagger(t,0)
\right].
\label{eq:phys-oqs-exact-reduced-map}
\end{equation}
该映射自动保持迹、厄米性和完全正性。为看见完全正性从哪里来，把环境初态谱分解为 $\rho_B=\sum_r p_r|r\rangle\langle r|$，并在末态环境基 $\{|\alpha\rangle\}$ 上取偏迹。定义
\begin{equation}
K_{\alpha r}(t)=\sqrt{p_r}\,{}_B\!\langle\alpha|U_{\rm tot}(t,0)|r\rangle_B,
\label{eq:oqs-kraus-operators}
\end{equation}
便得到 Kraus 表示
\begin{equation}
\Phi_t[\rho]=\sum_{\alpha,r}K_{\alpha r}(t)\rho K_{\alpha r}^\dagger(t).
\label{eq:oqs-kraus-map}
\end{equation}
总演化的幺正性给出
\[
\sum_{\alpha,r}K_{\alpha r}^\dagger K_{\alpha r}
=\sum_r p_r,{}_B\!\langle r|U_{\rm tot}^\dagger U_{\rm tot}|r\rangle_B=I_S,
\]
所以式 \eqnrefs{eq:oqs-kraus-map} 保持迹。若再给系统附加一个不参与演化的旁观参考系，映射仍保持联合正性；这正是“完全正”比“正”多出的要求。

Kraus 表示还可以直接变成一次完整计算。考虑二能级原子的振幅阻尼：原子激发态把一个量子放进初始真空环境，幺正演化在相关二维子空间中取为
\[
\begin{aligned}
U\lvert0\rangle_S\lvert0\rangle_B
&=\lvert0\rangle_S\lvert0\rangle_B,\\
U\lvert1\rangle_S\lvert0\rangle_B
&=\sqrt{1-p}\,\lvert1\rangle_S\lvert0\rangle_B
+\sqrt p\,\lvert0\rangle_S\lvert1\rangle_B,
\end{aligned}
\]
其中 \(0\leq p\leq1\) 是在所考察时间内发射量子的概率。环境初态只有
\(\lvert0\rangle_B\)，所以只需对末态环境基
\(\lvert0\rangle_B,\lvert1\rangle_B\) 取矩阵元：
\begin{equation}
K_0=\begin{pmatrix}1&0\\0&\sqrt{1-p}\end{pmatrix},
\qquad
K_1=\begin{pmatrix}0&\sqrt p\\0&0\end{pmatrix}.
\label{eq:oqs-amplitude-damping-kraus}
\end{equation}
逐项相乘可见
\(K_0^\dagger K_0+K_1^\dagger K_1=I\)。若
\[
\rho(0)=\begin{pmatrix}
\rho_{00}&\rho_{01}\\
\rho_{10}&\rho_{11}
\end{pmatrix},
\]
则
\begin{equation}
\rho(t)=
\begin{pmatrix}
\rho_{00}+p\rho_{11}&\sqrt{1-p}\,\rho_{01}\\
\sqrt{1-p}\,\rho_{10}&(1-p)\rho_{11}
\end{pmatrix}.
\label{eq:oqs-amplitude-damping-map}
\end{equation}
激发态人口以因子 \(1-p\) 减少，失去的人口全部流入基态；相干项则只乘 \(\sqrt{1-p}\)。迹保持条件负责人口守恒，而相干项与人口使用不同幂次，说明“衰减”不能用一个非幺正态矢量演化完整代替。取短时间
\(p=\gamma\,\Delta t+O(\Delta t^2)\)，把
\eqnrefs{eq:oqs-amplitude-damping-map}
减去初态再除以 \(\Delta t\)，就得到下一章会反复出现的自发辐射 Lindblad 耗散子。

固定同一个 $\rho_B$ 对所有系统初态使用，是上述 Kraus 形式的制备前提。若制备过程只允许系统初态落在集合 $\mathfrak D_{\mathcal A}$ 中，可以引入赋值映射
\begin{equation}
\mathcal A:\mathfrak D_{\mathcal A}\to\mathcal S(\Hil_S\otimes\Hil_B),
\qquad
\Tr_B\mathcal A[\rho]=\rho,
\label{eq:phys-oqs-assignment-map}
\end{equation}
并定义
\begin{equation}
\Phi_{t,0}^{\mathcal A}[\rho]
=\Tr_B\!\left[U_{\rm tot}(t,0)\mathcal A[\rho]U_{\rm tot}^\dagger(t,0)\right],
\qquad \rho\in\mathfrak D_{\mathcal A}.
\label{eq:phys-oqs-assignment-reduced-map}
\end{equation}
乘积赋值 $\mathcal A_0[\rho]=\rho\otimes\rho_B$ 定义在整个系统态空间上，并直接给出\eqnrefs{eq:phys-oqs-exact-reduced-map} 的完全正迹保持映射。若初始总态含系统--环境相关，\eqnrefs{eq:phys-oqs-exact-reduced-state} 对该给定初态仍然精确，但“只给系统初态 \(\rho(0)\) 就唯一决定总初态”的赋值通常只在一个兼容子集上定义。此时可以在这个兼容定义域上讨论约化映射，却一般不能把它延拓成定义在整个系统态空间上的同一个完全正迹保持映射。因而是否采用因子化初态首先是制备与定义域问题，不是 Born 弱耦合截断本身。

\eqnrefs{eq:phys-oqs-exact-reduced-map} 通常不能仅用有限个系统算符写成闭合的时间局域微分方程。\secref{sec:phys-oqs-gksl}将从它对应的相互作用绘景方程出发，依次引入 Born、马尔可夫与久期条件，构造可计算的 GKSL 生成元。


精确约化映射通常依赖系统的整个历史，不能只用当前的 $\rho(t)$ 闭合。下一章将从相互作用绘景的精确积分方程出发，逐步比较环境相关时间、系统驰豫时间和不同 Bohr 频率的分离；只有这些尺度条件成立，时间局域的 GKSL 方程才会出现。
